\section{Attribute Representation in CUB-200-2011}

CUB-200-2011 provides rich visual attribute annotations in addition to species labels. The dataset contains 312 bird attributes that describe visual properties such as color, shape, size, and pattern \cite{wah2011cub}. These attributes are especially useful for fine-grained bird recognition because many bird species differ only in subtle local traits, such as bill shape, crown color, wing pattern, belly color, or tail shape.

In this thesis, the CUB attributes are used as auxiliary supervision rather than as additional inputs during inference. Given an image, the model still predicts the taxonomic labels directly from the image. Attribute annotations are used only during training to guide the model toward discriminative bird traits.

\subsection{Attribute Groups}

In this thesis, the 312 CUB attributes are organized into 28 semantic groups based on their attribute names. Each group corresponds to a bird trait type, and attributes within the same group represent related states of the same visual property. For example, the group \textit{bill shape} contains different bill-shape states, while the group \textit{wing color} contains possible color states for the wing region. This group structure is important because attributes within the same group are mutually related and describe the same semantic property.

The attribute groups used in this thesis are shown in Table~\ref{tab:cub_attribute_groups}. The number of attribute states differs across groups. Color-related groups usually contain 15 states, while pattern-related groups usually contain fewer states.

\begin{table}[t]
    \centering
    \caption{CUB attribute groups used in this thesis. Each group contains a set of attribute states.}
    \label{tab:cub_attribute_groups}
    \begin{tabular}{l c l c}
        \hline
        \textbf{Attribute group} & \textbf{\# states} & \textbf{Attribute group} & \textbf{\# states} \\
        \hline
        bill shape & 9 & wing color & 15 \\
        upperparts color & 15 & underparts color & 15 \\
        breast pattern & 4 & back color & 15 \\
        tail shape & 6 & upper tail color & 15 \\
        head pattern & 11 & breast color & 15 \\
        throat color & 15 & eye color & 14 \\
        bill length & 3 & forehead color & 15 \\
        under tail color & 15 & nape color & 15 \\
        belly color & 15 & wing shape & 5 \\
        size & 5 & shape & 14 \\
        back pattern & 4 & tail pattern & 4 \\
        belly pattern & 4 & primary color & 15 \\
        leg color & 15 & bill color & 15 \\
        crown color & 15 & wing pattern & 4 \\
        \hline
    \end{tabular}
\end{table}

Let the set of all attribute groups be denoted as
\begin{equation}
    \mathcal{G} = \{g_1, g_2, \ldots, g_K\},
\end{equation}
where \(K=28\) for the full CUB attribute set. Each group \(g\) contains a set of attribute indices \(I_g\). The full 312-dimensional attribute vector for a species \(c\) is denoted as
\begin{equation}
    \mathbf{a}_c \in \mathbb{R}^{312}.
\end{equation}

The class-level attribute vector is taken from the CUB class-level continuous attribute annotation file, where each entry indicates how frequently the attribute is considered present for a species.

For an image \(x_i\) with species label \(y_i^s\), the corresponding attribute target is obtained from the class-level attribute annotation of that species:
\begin{equation}
    \mathbf{a}_i = \mathbf{a}_{y_i^s}.
\end{equation}



\subsection{Group-wise Attribute Prediction}

A straightforward design would be to predict the full 312-dimensional attribute vector directly. However, this thesis does not use that formulation. Instead, the model treats attribute prediction as a group-wise classification problem. This design follows the semantic organization of CUB attributes: attributes within the same group represent related states of the same visual property.

Although the original CUB annotations are provided as attribute-level presence labels, this thesis converts each group into a single-label classification target by selecting the highest-valued attribute state within the group.

For each attribute group \(g\), the model predicts one attribute state from the states belonging to that group. The target for group \(g\) is obtained by selecting the attribute state with the highest value inside the group:
\begin{equation}
    t_{i,g}
    =
    \arg\max_{j \in I_g} a_{i,j}.
\end{equation}
Here, \(t_{i,g}\) is the group-wise attribute target for sample \(i\) and group \(g\).

The model uses a separate classifier for each active attribute group. Let \(\mathbf{h}_{i,g}^{a}\) denote the attribute-oriented feature for group \(g\). The group-specific classifier produces logits
\begin{equation}
    \mathbf{z}_{i,g}^{a}
    =
    W_g \mathbf{h}_{i,g}^{a} + b_g,
\end{equation}
where \(W_g\) and \(b_g\) are the parameters of the classifier for group \(g\). The dimension of \(\mathbf{z}_{i,g}^{a}\) is equal to the number of retained attribute states in that group.

This means that the model does not predict all CUB attributes using a single shared binary classifier. Instead, it predicts attribute states through multiple group-specific classifiers. This formulation is more consistent with the structure of CUB attributes because each classifier specializes in one semantic trait type, such as bill shape, wing color, or tail pattern.

\subsection{Attribute Masks and Active Groups}

When all 312 attributes are used, each group classifier predicts among all states in its corresponding group. However, this thesis also studies selected attribute subsets, such as HSDAS-199, Sibling-IVAS-188, and Random-199. In these settings, only a subset of attributes is retained.

Let \(M_g \subseteq I_g\) denote the selected attribute indices for group \(g\). If \(M_g\) is non-empty, then \(g\) is considered an active attribute group. The active group set is defined as
\begin{equation}
    \mathcal{G}_{active}
    =
    \{g \in \mathcal{G} \mid |M_g| > 0\}.
\end{equation}

For each active group, the original attribute indices are remapped to local class indices. If the target attribute state of a sample is retained in the selected mask, it is used as the target for the group-specific attribute loss. If the target attribute state is not retained, the attribute target for that sample and group is ignored during loss computation. This avoids forcing the model to predict an attribute state that is not included in the selected subset.

Formally, for a selected group \(g\), the remapped target is
\begin{equation}
    \tilde{t}_{i,g}
    =
    \phi_g(t_{i,g}),
\end{equation}
where \(\phi_g(\cdot)\) maps the original CUB attribute index to the local index inside \(M_g\). If \(t_{i,g} \notin M_g\), then \(\tilde{t}_{i,g}\) is set to an ignore label and excluded from the attribute loss.

\subsection{Role in the Proposed Model}

The group-wise attribute representation serves two roles in the proposed method. First, it provides auxiliary supervision through group-specific attribute prediction heads. These heads encourage the model to learn features related to semantic bird traits. Second, the hidden attribute-oriented features are used for species-level fusion. Instead of using the predicted discrete attribute labels as inputs, the model uses continuous attribute-oriented representations learned before the group classifiers.

This distinction is important. The proposed model does not classify a bird by first predicting a complete attribute vector and then using that vector as the final representation. Rather, attribute prediction acts as an auxiliary task, while the learned attribute-oriented features are fused with the species-level visual representation. Therefore, the attribute module guides feature learning during training and supports species classification, while the inference process remains image-based.

Overall, this representation allows the model to exploit the semantic structure of CUB attributes without requiring manual attribute inputs at test time. It also enables different attribute selection strategies to be implemented through group masks, making it possible to compare FullAttr-312, HSDAS-199, Sibling-IVAS-188 and Random-199 under the same attribute-guided H-CAST framework.